\section{Ordinary averaging does not remove the canonical obstruction (NS-65)}
\label{sec:ns65-cesaro}

The preceding results leave different coefficient rules open. One simple
repair of the sharp canonical rule can be tested exactly: ordinary
Ces\`aro averaging over its index. The conclusion below concerns a fixed
number of these averages. Logarithmic tapers, a growing number of averages,
and optimized coefficients are outside its scope.

For a sequence $(f_N)$ of vectors or scalars define
\[
 (Cf)_N=\frac1N\sum_{n=1}^Nf_n,\qquad
 B_N^{[q]}=(C^qB)_N,\qquad
 H_N^{[q]}=(C^q a)_N,\quad a_n=n\eta_n,
\]
where $B_n$ and $\eta_n$ are the canonical objects of NS-62 and
$q\geq0$ is a fixed integer. The averaging weights are nonnegative,
sum to one, and only involve indices $n\leq N$.

\begin{proposition}[A fixed number of Ces\`aro averages leaves a critical-zero pole]
\label{prop:ns65-cesaro-obstruction}
For every fixed integer $q\geq0$ and every sequence of nonnegative
integers $r_N$,
\[
 \norm{J^{r_N}(\chi+B_N^{[q]})}_H\not\longrightarrow0
 \quad(N\to\infty\text{ through all integers}).
\]
More precisely, for every $N$ and every $r\geq0$,
\begin{equation}
 \norm{J^r(\chi+B_N^{[q]})}_H
 \geq\frac{|H_N^{[q]}-k_{N,r}|}{\sqrt N}
 \geq\frac{|H_N^{[q]}-2|}{\sqrt N}
             -\frac{4\pi^2}{3^{r+2}N^{5/2}},
 \label{eq:ns65-averaged-witness}
\end{equation}
and $H_N^{[q]}\ne o(\sqrt N)$. This is full-sequence nonconvergence;
no selected-subsequence exclusion or assertion about signed optimal
approximation follows.
\end{proposition}
\begin{proof}
For $0<x<1/N$, every index $n\leq N$ obeys the same small-cell
identity from NS-62. Linearity and normalization of the weights give
\[
 \mathcal V(\chi+B_N^{[q]})(x)=k(x)-H_N^{[q]}.
\]
Apply the same unit Laguerre test $\psi_{N,r}$ and its complete
adjoint identity. This proves~\eqref{eq:ns65-averaged-witness},
including its remainder uniform in $r$.

It remains to prove the scalar obstruction. Put
\[
 \eta(x)=\int_1^x\frac{M(t)}{t^2}\,dt,\qquad
 A_0(x)=x\eta(x),\qquad
 A_{q+1}(x)=\frac1x\int_1^x A_q(u)\,du.
\]
For $\Re s>1$, absolute convergence and Fubini give
\begin{equation}
 \int_1^\infty A_q(x)x^{-s-1}\,dx
   =\frac1{(s+1)^q(s-1)s\zeta(s)}.
 \label{eq:ns65-cesaro-mellin}
\end{equation}
For $q=0$ this is the Mellin identity for $\eta$ with its variable
shifted by one. Each continuous averaging contributes exactly the
factor $(s+1)^{-1}$; there is no discarded endpoint term.

We now connect this continuous identity to the discrete averages.
Directly $A_0(x)=O(x\log(2x))$ and, almost everywhere,
\[
 A_0'(x)=\eta(x)+M(x)/x
         =\sum_{n\leq x}\mu(n)/n=O(\log(2x)).
\]
Inductively $A_q(x)=O_q(x\log(2x))$ and
$A_{q+1}'(x)=(A_q(x)-A_{q+1}(x))/x=O_q(\log(2x))$.
The functions are locally absolutely continuous. Comparing each
unit-cell integral with its endpoint value therefore gives
\[
 \frac1N\sum_{n\leq N}A_q(n)-\frac1N\int_1^N A_q(x)\,dx
       =O_q(\log(2N)).
\]
Since $A_0(N)=a_N$, induction yields
\begin{equation}
 H_N^{[q]}=A_q(N)+O_q(\log(2N)).
 \label{eq:ns65-discrete-continuous}
\end{equation}
Thus $H_N^{[q]}=o(\sqrt N)$ would imply $A_q(x)=o(\sqrt x)$
also between integers, by the displayed derivative bounds.

Under that hypothetical little-$o$ estimate the Mellin integral in
\eqref{eq:ns65-cesaro-mellin} is analytic for $\Re s>1/2$.
At a known critical-line zero $\rho=1/2+i\gamma_*$, its usual
Abelian boundary estimate gives
\[
 \lim_{\sigma\downarrow1/2}(\sigma-1/2)
       \int_1^\infty A_q(x)x^{-\sigma-i\gamma_*-1}\,dx=0.
\]
For example, split the integral at a fixed large $X$, use
$|A_q(x)|\leq\varepsilon\sqrt x$ on the tail, then let
$\sigma\downarrow1/2$ and $\varepsilon\downarrow0$.
But the meromorphic continuation on the right side of
\eqref{eq:ns65-cesaro-mellin} has a genuine pole at $\rho$:
none of $\rho+1$, $\rho-1$ or $\rho$ is zero, and its numerator
is one. A pole of any positive order contradicts the boundary limit.
Therefore $\limsup_N|H_N^{[q]}|/\sqrt N>0$.
Together with~\eqref{eq:ns65-averaged-witness}, whose remainder tends
to zero uniformly in $r\geq0$, this proves the stated nonconvergence.
\end{proof}

\paragraph{Interpretation and stop condition.}
Ordinary averaging attenuates a critical-zero contribution by
$(\rho+1)^{-q}$, but a fixed $q$ cannot annihilate it. This is an
analytic obstruction for the specified averaged sharp rule, not a
numerical observation or a priority claim about summability methods.
The general summability viewpoint is already part of
\cite[Section 3]{BaezDuarteStrong}. The useful open task remains an
independent estimate for the optimized signed residual, or a genuinely
different coefficient rule with a complete norm bound. No such lower
correlation or convergence estimate is supplied here.
